\documentclass[12pt]{oci}
\usepackage{graphicx}
\usepackage{enumitem}
\usepackage{titling}

\phase{Clasificatoria Regional}
\problemcnttrue%

\begin{document}

\begin{titlingpage}
  \begin{center}
  \includegraphics[scale=0.6]{logo.eps}

  \vskip 70pt
  \Large{\bf Olimpiada Chilena de Inform\'atica\\ \the\year}
  \vskip 10pt
  \large{Clasificatoria Regional}
  \vskip 10pt
  \normalsize{\it 11 de Octubre, \the\year}


  \vskip 95pt

  \end{center}
  Este conjunto de problemas est\'a siendo utilizado de manera simultanea en las siguientes sedes:

  \begin{itemize}[leftmargin=*]
  \item Arica
  \item Santiago
  \item Valparaiso
  \item Curic\'o
  \item Talca
  \item Chill\'an
  \item Concepci\'on
  \item Temuco
  \end{itemize}
\end{titlingpage}


\subsection*{Informaci\'on General}

Esta p\'agina muestra informaci\'on general que se aplica a todos los problemas.

\subsection*{Env\'io de una Soluci\'on}


\begin{enumerate}
\itemsep 0em
\item Los participantes deben enviar {\bf un solo archivo} con el c\'odigo fuente de su soluci\'on.
\item El nombre del archivo debe tener la extensi\'on \verb+.c+, \verb+.cpp+ o \verb+.pas+ dependiendo de si la soluci\'on est\'a escrita en \verb|C|, \verb|C++| o \verb|Pascal| respectivamente.
\end{enumerate}

\subsection*{Casos de Prueba y Subtareas}
\begin{enumerate}
\itemsep 0em
\item La soluci\'on enviada por los participantes ser\'a ejecutada varias veces con distintos casos de prueba.
\item Cada problema define diferentes subtareas que restringen el problema y se asignar\'a puntaje de acuerdo a la cantidad de subtareas que logre solucionar de manera correcta.
\item Una soluci\'on puede resolver al mismo tiempo m\'as de una subtarea.
\item La soluci\'on es ejecutada con cada caso de prueba de manera independiente y por tanto puede fallar en algunas subtareas sin influir en la ejecuci\'on de otras.
\end{enumerate}

\subsection*{Entrada}
\begin{enumerate}
\itemsep 0em
\item Toda lectura debe ser hecha desde la {\bf entrada est\'andar} usando, por ejemplo, las funciones \verb+scanf+, \verb+std::cin+ o \verb+readln+ dependiendo del lenguaje escogido.
\item La entrada corresponde a un solo caso de prueba el cual est\'a descrito en varias l\'ineas dependiendo del problema.
\item {\bf Se debe asumir que la entrada sigue el formato descrito} en el enunciado de cada problema.
\end{enumerate}

\subsection*{Salida}
\begin{enumerate}
\itemsep 0em
\item Toda escritura debe ser hecha hacia la {\bf salida est\'andar} usando, por ejemplo, las funciones \verb+printf+, \verb+std::cout+ o \verb+writeln+ dependiendo del lenguaje escogido.
\item El resultado del caso de prueba debe ser escrito siguiendo el formato descrito para cada problema.
\item {\bf El formato de salida debe ser seguido de manera estricta} considerando los espacios especificados, las may\'usculas y min\'usculas.
\item Toda l\'inea, incluyendo la \'ultima, debe terminar con un salto de l\'inea.
\end{enumerate}


\end{document}
